\documentclass[11pt,a4paper]{article}

\usepackage[utf8]{inputenc}
\usepackage[T1]{fontenc}
\usepackage[spanish,es-noquoting]{babel}
\usepackage{amsmath,amssymb}
\usepackage{graphicx}
\usepackage{siunitx}
\usepackage{booktabs}
\usepackage[hidelinks]{hyperref}
\usepackage[margin=2.5cm]{geometry}

\title{Búsqueda del bosón $Z$ en datos abiertos de ATLAS}
\author{Andrés}
\date{\today}

\begin{document}

\maketitle

\begin{abstract}
Resumen del análisis: reconstrucción del bosón $Z$ a partir del canal
dileptónico usando los datos abiertos de ATLAS.
\end{abstract}

\section{Introducción}

El bosón $Z$ es el mediador neutro de la interacción débil. Tiene carga
eléctrica nula y una masa de aproximadamente \SI{91.19}{\giga\electronvolt}.

\section{Método}

Se reconstruye la masa invariante del par de leptones:

\begin{equation}
    m_{\ell\ell} = \sqrt{2\, p_{T,1}\, p_{T,2}
    \left[\cosh(\eta_1 - \eta_2) - \cos(\phi_1 - \phi_2)\right]}
\end{equation}

\section{Resultados}

% Para insertar una figura exportada desde el notebook:
% \begin{figure}[htbp]
%     \centering
%     \includegraphics[width=0.8\textwidth]{figuras/masa_invariante.pdf}
%     \caption{Distribución de masa invariante dileptónica.}
%     \label{fig:mll}
% \end{figure}

\section{Conclusiones}

\end{document}
